%% P03 --- Rzędy wielkości: notacja O, FLOP-y i pamięć
%%
%% Zalążek, wygenerowany przez tools/gen_stubs.py z tools/programs.json.
%% Poniższy opis jest kontraktem tego programu: co ma uzasadnić i co ma dać
%% czytelnikowi. Napisanie programu oznacza usunięcie bloku \programstub{} i
%% zastąpienie go programem -- po czym ten plik nie jest już generowany.

\program{Rzędy wielkości: notacja O, FLOP-y i pamięć}
\label{prog:P03}

\programstub{%
\textit{[Opis do przetłumaczenia --- patrz wydanie angielskie.]}

Argument: asymptotic notation is a statement about growth and says nothing
about which of two implementations is faster today; both facts matter.
Payoff: bytes, FLOPs as a rate, and arithmetic intensity as a ratio, which
is the whole justification for kernel fusion; why a large matrix multiply is
compute-bound and an elementwise operation is memory-bound. DEPENDENCY NOTE:
an earlier brief had the reader count the parameters of a transformer block
here, three parts before any program has said what a matrix is. That count
moves to P32, where the shapes exist. What stays is magnitude and cost,
which need only arithmetic and sigma notation.

\medskip
\textbf{Planowana długość:} 45 ramek.
}
